\documentclass[tikz,border=8pt]{standalone}
\usepackage{fontspec}
\usepackage{xeCJK}
\IfFontExistsTF{Segoe UI}{\setmainfont{Segoe UI}}{\setmainfont{Arial}}
\IfFontExistsTF{Microsoft JhengHei}{\setCJKmainfont{Microsoft JhengHei}}{\IfFontExistsTF{Noto Sans CJK TC}{\setCJKmainfont{Noto Sans CJK TC}}{\setCJKmainfont{Noto Sans TC}}}
\IfFontExistsTF{Microsoft JhengHei}{\setCJKsansfont{Microsoft JhengHei}}{\IfFontExistsTF{Noto Sans CJK TC}{\setCJKsansfont{Noto Sans CJK TC}}{\setCJKsansfont{Noto Sans TC}}}
\xeCJKsetup{CJKglue={\hskip 0.045em plus 0.015em minus 0.005em}, CJKecglue={\hskip 0.08em plus 0.02em}}
\IfFileExists{../../../FontAwesome.otf}{\newfontfamily\FA[Path=../../../]{FontAwesome.otf}}{\newfontfamily\FA{Segoe UI Symbol}}
\newcommand{\faIcon}[1]{{\FA\symbol{#1}}}
\usetikzlibrary{arrows.meta,positioning,calc}
\definecolor{paper}{HTML}{FFFDF8}
\definecolor{navy}{HTML}{0F172A}
\definecolor{muted}{HTML}{334155}
\definecolor{line}{HTML}{CBD5E1}
\definecolor{blue}{HTML}{2563EB}
\definecolor{bluesoft}{HTML}{DBEAFE}
\definecolor{green}{HTML}{00856A}
\definecolor{greensoft}{HTML}{DCFCE7}
\definecolor{red}{HTML}{DC2626}
\definecolor{redsoft}{HTML}{FEE2E2}
\begin{document}
\begin{tikzpicture}[x=1cm,y=1cm,>=Latex,line cap=round,line join=round]
\path[fill=paper] (0,0) rectangle (16,9.6);
\node[font=\bfseries\fontsize{18}{23}\selectfont,text=navy,align=center,text width=14.5cm] at (8,8.82) {最後一輪必須回到人的剪輯桌};
\node[font=\fontsize{10.3}{13.8}\selectfont,text=muted,align=center,text width=13.2cm] at (8,7.94) {LLM 很會補齊；作者的工作，是刪掉沒有判斷的順滑句子。};

\draw[rounded corners=12pt,fill=bluesoft,draw=blue,line width=1pt] (1.0,4.35) rectangle (4.3,6.15);
\node[font=\bfseries\fontsize{11}{14}\selectfont,text=navy,align=center,text width=2.75cm] at (2.65,5.62) {模型初稿};
\node[font=\fontsize{8.2}{10.2}\selectfont,text=muted,align=center,text width=2.7cm] at (2.65,5.12) {完整、平均\\沒有風險};

\draw[rounded corners=12pt,fill=white,draw=line,line width=1pt] (6.15,2.45) rectangle (9.85,7.0);
\node[font=\bfseries\fontsize{12}{15}\selectfont,text=navy,align=center,text width=3.2cm] at (8,6.35) {人的剪輯};
\draw[rounded corners=8pt,fill=redsoft,draw=red,line width=.9pt] (6.75,5.05) rectangle (9.25,5.75);
\node[font=\fontsize{8.4}{10.4}\selectfont,text=navy,align=center,text width=2.25cm] at (8,5.4) {刪掉空話};
\draw[rounded corners=8pt,fill=greensoft,draw=green,line width=.9pt] (6.75,4.05) rectangle (9.25,4.75);
\node[font=\fontsize{8.4}{10.4}\selectfont,text=navy,align=center,text width=2.25cm] at (8,4.4) {補上證據};
\draw[rounded corners=8pt,fill=white,draw=line,line width=.9pt] (6.75,3.05) rectangle (9.25,3.75);
\node[font=\fontsize{8.4}{10.4}\selectfont,text=navy,align=center,text width=2.25cm] at (8,3.4) {留下立場};

\draw[rounded corners=12pt,fill=navy,draw=navy,line width=1pt] (11.7,4.35) rectangle (15.0,6.15);
\node[font=\bfseries\fontsize{10.4}{13}\selectfont,text=white,align=center,text width=2.6cm] at (13.35,5.68) {可以署名\\的作品};
\node[font=\fontsize{8.0}{10}\selectfont,text=white,align=center,text width=2.55cm] at (13.35,4.95) {有取捨\\也有責任};

\draw[->,draw=navy,line width=1pt] (4.42,5.25)--(6.0,5.25);
\draw[->,draw=navy,line width=1pt] (10.0,5.25)--(11.55,5.25);
\draw[draw=line,line width=.9pt] (1.0,1.55)--(15.0,1.55);
\node[font=\fontsize{10.2}{12.8}\selectfont,text=navy,align=center,text width=12.6cm] at (8,.95) {如果最後一版看不出作者願意負責的句子，那只是乾淨，不是完成。};
\end{tikzpicture}
\end{document}
